\subsection{问题一的分析}

问题一的核心任务是基于附件1和附件2的数据，精确计算冷链药品运输的总费用。该问题的关键在于：（1）从排货信息中准确提取目的地城市，并将其映射到地理坐标以计算运输距离；（2）建立符合实际运营逻辑的运输成本模型。基础的直线距离计算可采用Haversine球面距离公式，但实际道路距离需要考虑绕行因素，且不同地区（东部、中部、西部）的道路网络密度差异较大，应采用区域差异化的绕行系数。成本模型需要区分可变成本（油费、路桥费）和固定成本（人工、保险、折旧等），在往返运输场景下还要考虑回程空载带来的油耗下降。问题一属于数据驱动的成本核算问题，求解的关键在于数据清洗的鲁棒性和成本分摊逻辑的合理性。

\subsection{问题二的分析}

问题二要求在问题一的成本模型基础上，通过拼车合并运输实现成本节约。这是一个带容量约束的车辆路径问题（Capacitated Vehicle Routing Problem, CVRP）的变体：从单一仓库出发，需要将分布在不同城市的订单合并为尽可能少的运输路线。与标准CVRP的区别在于：（1）成本函数是非线性的——车辆选择（4.2M/7.6M/9.6M三种冷藏车）影响单位成本和载货容量，需匹配实际货物托盘数；（2）存在方向可行性约束——只有方向相近的目的地才适合拼车，否则绕行成本将抵消节约；（3）存在时间窗口约束——每日装卸时间仅8小时（9:00–17:00），每次装卸2小时。考虑到数据中存在自然的"周"分组结构，可逐周进行优化，并探索相邻周之间的跨周合并空间。Clarke-Wright节约算法因其计算效率高、易于施加方向约束的特性，适合本问题的求解。TSP路径排序（最近邻构造+2-opt改进）可用于优化同一路线内的访问顺序。

\subsection{问题三的分析}

问题三的性质与问题一、二不同——它不直接求解一个优化问题，而是先评估现有方案的合理性，再提出改进。数据规模更大（一周约53.6万条运单、22.2万条派车记录），涉及152个城市、88个仓库。评估的核心挑战在于：（1）设计合理的评估指标体系——不能主观判断，需要从配送效率、路线质量、时效表现、需求结构、合规水平等多个维度综合度量；（2）权重的确定——不同指标的重要程度不同，且应避免主观赋权带来的偏差。熵权法是一种基于数据本身差异程度的客观赋权方法：指标在各仓库间的差异越大，说明该指标的区分能力越强，应赋予更高权重。优化方面，派车单安排存在三层改进空间：单车辆日内的TSP路径排序→同一仓库同一天内不同车辆间的订单合并→合并后路线与最优车型的匹配。

\subsection{流程分析}

整体求解流程如下：首先进行数据预处理，包括城市名提取、地理坐标匹配、运输趟次/车辆日聚合；然后分别针对三个问题建立模型——问题一的改进运输成本核算模型、问题二的Clarke-Wright拼车优化模型、问题三的熵权法合理性评估模型与三层递进优化模型；接着求解并输出数值结果；最后进行模型评价。三问题之间存在递进关系：问题一提供基础成本核算框架，问题二在此基础上引入拼车优化，问题三则将评估与优化拓展到多仓库、多终端的大规模实际运作场景。
